% Core formal section for
% "What an Extensional First-Order Formalization Leaves Underdetermined:
%  Stoic Katalepsis and Humean Custom"
%
% Revision: 2026-08-17
%   - P0: thesis narrowed to "the target relation is not determined by the
%         L_0-reduct over the admissible class K";
%   - P0: Proposition 1 (formerly "Theorem 1") explicitly conditional;
%   - P0: Bridge-collapse proposition and characterization added;
%   - P0: Proposition 2 rebuilt as an implicit-definability failure,
%         with definability, stability and Obtains sub-sort;
%   - P0: Hume reading labelled "Strong exclusion reconstruction";
%   - P0: Stoic formula relabelled "minimal dependency skeleton";
%   - P1: language-level discipline (FOL, metalanguage, modal, JL);
%   - P1: separate Cont / Fact sorts with Obtains monad;
%   - P1: modal operator \Box_s indexed explicitly;
%   - P1: provenance table updated; comparative section deferred to appendix.
%
% Intended for integration into the main manuscript as \S2 of §2.

\subsection{Interpretive and Formal Preconditions}
\label{sec:formal-preconditions}

This is a methodological note, not a historical thesis. Its aim is to
state precisely what an extensional first-order encoding at a fixed
level of granularity can and cannot represent, and to keep the resulting
claims conditional on the encoding choices that produced them. We do
not attribute a contemporary grounding relation, a fully articulated
theory of epistemic justification, or a complete process-reliabilistic
framework to the historical authors. Where we use such vocabulary, it
is a contemporary regimentation of one strong interpretive option; an
alternative reading is recorded alongside it.

The formal reconstruction is deliberately narrower than a reconstruction
of either Stoic epistemology or Hume's epistemology as a whole. Its
purpose is to isolate the minimum content-bearing structure required for
the representational argument. Accordingly, we bracket the full Stoic
ontology of \emph{lekta} and retain only the structure needed to
distinguish impressions, represented contents, cognitive episodes,
belief-production, and justification. In particular, the predicates
introduced below do not exhaust, nor do they uniquely capture, the
historical vocabularies.

\begin{quote}
For purposes of reconstruction, we bracket the full Stoic ontology of
\emph{lekta} and Hume's full apparatus of passions, sympathy and
moral psychology, and represent only the minimal content-bearing
structure required for the formal result. Where an interpretive
choice is required, the choice is stated explicitly and an
alternative reading is acknowledged.
\end{quote}

\subsection{The Minimal Extensional Language $L_0$}
\label{sec:l0}

Let $L_0$ be a many-sorted first-order language with the following
core sorts:
\[
I \quad\text{(impressions)},\qquad
\mathit{Cont} \quad\text{(ordinary representational contents)},\qquad
B \quad\text{(cognitive episodes)}.
\]
A subsort $\mathit{Fact} \subseteq \mathit{Cont}$ is introduced only in
the enriched language $L^+$ of \S\ref{sec:enrichment} to host
reified grounding \emph{relata}; it is not part of $L_0$ proper.
An additional sort $E$ for explicit justification terms is introduced
only in the justification-logic enrichment of
\S\ref{sec:richer-languages} and is not part of the core extensional
language.

The non-logical vocabulary of $L_0$ is:

\begin{align*}
&\operatorname{Kat}(i)
&&\text{: $i\in I$ is kataleptic;} \\
&\operatorname{Rep}(i,c)
&&\text{: impression $i$ represents content $c\in\mathit{Cont}$;} \\
&\operatorname{Grasp}(b,i)
&&\text{: cognitive episode $b$ grasps impression $i$;} \\
&\operatorname{Assent}(b,c)
&&\text{: $b$ assents to content $c$;} \\
&\operatorname{Bel}(b,c)
&&\text{: $b$ is a belief state with content $c$ (whether or not
causally formed);} \\
&\operatorname{Causal}(b,c)
&&\text{: $b$ is a causal belief with content $c$;} \\
&\operatorname{Custom}(b)
&&\text{: $b$ is produced by custom;} \\
&\operatorname{Just}(b,c)
&&\text{: the relation of $b$ to $c$ is rationally
justified;}\footnotemark \\
&\operatorname{StoicEp}(b)
&&\text{: $b$ is a cognitive episode to which the reconstructed
Stoic condition applies.}
\end{align*}
\footnotetext{$\operatorname{Just}(b,c)$ is binary rather than unary. The
claim that a cognitive episode is ``justified'' \emph{simpliciter} is
distinct from the claim that its assent to a particular content is
justified. The choice of $b$ as ``episode'' (rather than
``assent-token'' or ``agent-relative state'') is itself a hermeneutic
decision, recorded here so that no later inference silently presupposes
a particular reading of Hume's ``belief''.}

This separation is intentional. In particular,
$\operatorname{Grasp}(b,i)$ is not conflated with representation
$\operatorname{Rep}(i,c)$ or assent $\operatorname{Assent}(b,c)$:
$\operatorname{Grasp}$ operates in the impression-impression space,
$\operatorname{Assent}$ in the content space, and $\operatorname{Bel}$
in the cognitive-state space. The distinction prevents later
inferences from silently collapsing what an extensional encoding
should keep separate.

\subsubsection{The Stoic Minimal Dependency Skeleton}
\label{sec:stoic-skeleton}

The Stoic formula is presented here as a \emph{minimal dependency
skeleton}: a relational pattern extracted from the katalepsis
discussion that is sufficient to register the relevant dependency
between kataleptic impression, representation, grasping and assent,
without claiming to exhaust Stoic epistemology. Modal and doxastic
content is recorded separately (cf.\ \S\ref{sec:stoic-modal-note}).

\[
\forall b\,
\Bigl(
\operatorname{StoicEp}(b)
\rightarrow
\exists i\,\exists c\,
\bigl[
\operatorname{Kat}(i)\land
\operatorname{Rep}(i,c)\land
\operatorname{Grasp}(b,i)\land
\operatorname{Assent}(b,c)
\bigr]
\Bigr).
\tag{S}
\]

Formula (S) is a reconstruction constraint, not a claim that the
complete Stoic epistemology reduces to this formula; in particular,
(S) does not by itself guarantee that $b$ is a belief state about
$c$, that the \emph{content} $c$ figured in both $\operatorname{Rep}$
and $\operatorname{Assent}$ is the same content, or that
$\operatorname{Kat}(i)$ is connected to a truth or ``could not arise
from what is not'' clause. The first of these is a deliberate choice:
the link from the minimal skeleton to genuine belief is the place
where the modern interpreter is owed a separate, content-level
argument, and we do not pre-empt that argument here.

\subsubsection{Humean Mechanism and the Exclusion Readings}
\label{sec:hume-readings}

The core formalization records two mechanism-level constraints:

\[
\forall b\,\forall c\,
\bigl(\operatorname{Causal}(b,c)\rightarrow \operatorname{Custom}(b)\bigr),
\tag{M$_{0a}$}
\]
\[
\forall b\,\forall c\,
\bigl(\operatorname{Causal}(b,c)\rightarrow \operatorname{Bel}(b,c)\bigr).
\tag{M$_{0b}$}
\]

No axiom is imposed that identifies custom-production with the
\emph{absence} of justification. Instead, three readings are
distinguished, each as a candidate reconstruction.

\paragraph{Strong exclusion reconstruction (H-I).}
This is one contemporary regimentation of the strongest reading of
Hume's claim that custom-production \emph{excludes} rational
justification:
\[
\forall b\,\forall c\,
\bigl[
\operatorname{Custom}(b)\land
\operatorname{Bel}(b,c)
\rightarrow
\neg\operatorname{Just}(b,c)
\bigr].
\tag{H-I}
\]
H-I is a candidate, not a presupposition: a contemporary interpreter
who takes Hume's talk of ``custom'' to carry this strong normative
weight will write (H-I) into her reconstruction; one who does not, will
not. The choice is hermeneutic, not formal.

\paragraph{Naturalistic-compatible reading (H-II).}
The weaker \emph{naturalistic-compatible} reading imposes no such
axiom. On H-II, $\operatorname{Custom}(b)$ records a
belief-forming mechanism while the justificatory status of the
relation between $b$ and $c$ remains separately determined. The
distinction prevents the illicit inference from psychological genesis
to epistemic invalidity. We treat (H-II) as the default for the rest
of this section, in line with the methodological note's commitment
not to over-read Hume.

\paragraph{Direct-attachment reading ($T_1$).}
For comparison, define
\[
\forall b\,\forall c\,
\bigl[
\operatorname{Causal}(b,c)
\rightarrow
\neg\operatorname{Just}(b,c)
\bigr].
\tag{T$_1$}
\]
$T_1$ attaches the exclusion directly to causal beliefs; it is
logically stronger than H-I only when a further bridge axiom is
present (cf.\ Proposition~\ref{prop:bridge-collapse}).

The next result concerns the extensional relationship between
these formulations.

\subsection{Proposition~1: Strength Relation Between Two Reconstructions}
\label{sec:prop1}

\begin{proposition}[Strength relation between two proposed reconstructions]
\label{prop:strength}
Let $M_0$ be the conjunction of (M$_{0a}$) and (M$_{0b}$), let
$T_1$ be the formula above, and let $T_2$ denote (H-I). Then
\[
T_2\land M_0 \models T_1,
\]
whereas
\[
T_1\land M_0 \not\models T_2.
\]
\end{proposition}

\begin{proof}
For the first entailment, let $\mathcal M\models T_2\land M_0$ and
suppose $\mathcal M\models\operatorname{Causal}(b,c)$. By
(M$_{0a}$), $\mathcal M\models\operatorname{Custom}(b)$, and by
(M$_{0b}$), $\mathcal M\models\operatorname{Bel}(b,c)$. Hence the
antecedent of $T_2$ is true, and since $\mathcal M\models T_2$,
$\mathcal M\models\neg\operatorname{Just}(b,c)$. Therefore
$\mathcal M\models T_1$. Since $b,c$ were arbitrary,
$T_2\land M_0\models T_1$.

For the converse, consider a structure with one cognitive episode
$b_1$ and one content $c_1$ such that
$\operatorname{Causal}(b_1,c_1)=\mathbf{false}$,
$\operatorname{Custom}(b_1)=\mathbf{true}$,
$\operatorname{Bel}(b_1,c_1)=\mathbf{true}$,
$\operatorname{Just}(b_1,c_1)=\mathbf{true}$.
Then (M$_{0a}$) and (M$_{0b}$) are satisfied vacuously, $T_1$ is
satisfied vacuously, but $T_2$ is false (its antecedent is true and
its consequent is false). Hence $T_1\land M_0\not\models T_2$.
\end{proof}

Proposition~1 is purely comparative and conditional. It establishes
neither that Hume endorsed $T_2$, nor that the causal--custom relation
is correctly represented by $M_0$, nor that $T_1$ exhausts
Humean scepticism. The result is about two proposed reconstructions,
relative to a fixed signature, on a fixed class of models.

\subsubsection{Bridge collapse and characterization}
\label{sec:bridge-collapse}

A reader may object that the difference between $T_1$ and $T_2$ is
trivial once a natural interpretive decision is made. The
\emph{bridge axiom}
\[
\forall b\,\forall c\,
\bigl[
(\operatorname{Custom}(b)\land\operatorname{Bel}(b,c))
\rightarrow
\operatorname{Causal}(b,c)
\bigr]
\tag{B$_0$}
\]
encodes the obvious Humean gloss that every custom-produced belief
\emph{is} a causal belief (EHU 5.1; T 1.3.14.20). Under $B_0$, the two
readings collapse.

\begin{proposition}[Bridge collapse]
\label{prop:bridge-collapse}
Over $M_0$,
\[
T_1\land B_0 \models T_2.
\]
Given Proposition~\ref{prop:strength}, $T_1$ and $T_2$ are therefore
extensionally \emph{equivalent} over $M_0\land B_0$. Moreover, at a
fixed pair $(b,c)$ the two readings differ---i.e.\ $T_1$ holds at the
pair while $T_2$ fails there---if and only if
\[
\operatorname{Custom}(b)\land
\operatorname{Bel}(b,c)\land
\operatorname{Just}(b,c)\land
\neg\operatorname{Causal}(b,c).
\tag{$\star$}
\]
Consequently, over $M_0$, a model separates $T_1$ from $T_2$ (i.e.\
satisfies $T_1\land M_0\land\neg T_2$) if and only if it satisfies
$T_1\land M_0$ together with $\exists b\,\exists c\,(\star)$.
\end{proposition}

\begin{proof}
For the entailment, suppose $\mathcal M\models T_1\land M_0\land B_0$
and $\mathcal M\models\operatorname{Custom}(b)\land\operatorname{Bel}(b,c)$.
By $B_0$, $\mathcal M\models\operatorname{Causal}(b,c)$, and by $T_1$,
$\mathcal M\models\neg\operatorname{Just}(b,c)$. Hence $T_2$ holds.
For the characterization: at a fixed pair $(b,c)$, $T_1$ holds and
$T_2$ fails exactly when $(\star)$ holds there. Reading $(\star)$
existentially, $T_1\land M_0\land\neg T_2$ entails $(\star)$:
$\neg T_2$ yields a pair with $\operatorname{Custom}\land
\operatorname{Bel}\land\operatorname{Just}$, on which $T_1$ forces
$\neg\operatorname{Causal}$. Conversely, $T_1\land M_0$ together with
$(\star)$ entails $\neg T_2$ ($(\star)$ is a witness) and hence entails
$T_1\land M_0\land\neg T_2$; the two directions give the equivalence.
Note that $(\star)$ alone does not entail $T_1$: a model may satisfy
$(\star)$ at one pair while violating $T_1$ at another.
\end{proof}

The computational verification in \S\ref{sec:verification} confirms
both that $T_1\land M_0\land B_0\models T_2$ (zero falsifying
interpretations in the finite Boolean fragment) and that $(\star)$ is
the exact condition of separation at a single pair (globally,
separation is $T_1\land M_0$ together with $\exists b\,\exists c\,
(\star)$, as established above). The choice between $T_1$ and $T_2$
is therefore not settled by the logical facts alone; it is settled
by whether one accepts $B_0$. $B_0$ is an interpretive decision: it
identifies the predicate \texttt{Custom} with a mechanism
representable in the \texttt{Causal} vocabulary. That identification
is exactly the \emph{annotation gap} at the heart of this section.
The same kind of gap will reappear, in a more pointed form, in the
reification argument of \S\ref{sec:enrichment}.

\subsection{Controlled Reification and the Enriched Language}
\label{sec:enrichment}

The central difficulty arises when the intended Humean claim is not
merely that two facts co-occur but that one fact \emph{explains} or
\emph{grounds} another. Grounding is hyperintensional: it is not a
truth-functional connective, and it cannot be typed as a relation
between formulas in the object language. Accordingly, the
enrichment below uses \emph{controlled reification}: the explanatory
relata are first reified as fact-like objects, and grounding is then
typed as a relation between those objects. This is a
representational device, not a metaphysical commitment to grounding
facts as ordinary propositions of the historical theories.

Let the content sort $\mathit{Cont}$ also host, in the enriched
language, an additional subsort $\mathit{Fact}\subseteq\mathit{Cont}$.
A unary monadic predicate $\operatorname{Obtains}$ on $\mathit{Fact}$
marks the fact-like objects that actually obtain:
\[
\operatorname{Obtains}(f) \quad (f\in\mathit{Fact})
\]
is read ``$f$ obtains''. $\operatorname{Obtains}$ is part of $L^+$
but not of $L_0$.

Introduce two function symbols
\[
\operatorname{custFact}:B\to\mathit{Fact},
\qquad
\operatorname{nonjustFact}:B\times\mathit{Cont}\to\mathit{Fact},
\]
with the intended readings:
\begin{quote}
$\operatorname{custFact}(b)$ is the fact that $b$ is
custom-produced; $\operatorname{nonjustFact}(b,c)$ is the fact that
the relation of $b$ to $c$ lacks justification.
\end{quote}
The reification is disciplined by three bridge axioms, all in $L^+$:
\begin{align*}
\operatorname{Obtains}(\operatorname{custFact}(b))
&\leftrightarrow \operatorname{Custom}(b), \tag{O$_1$} \\
\operatorname{Obtains}(\operatorname{nonjustFact}(b,c))
&\leftrightarrow \neg\operatorname{Just}(b,c), \tag{O$_2$} \\
G(x,y) &\rightarrow \operatorname{Obtains}(x)\land\operatorname{Obtains}(y),
\tag{O$_3$}
\end{align*}
where $G$ is a new binary relation
\[
G\subseteq \mathit{Fact}\times\mathit{Fact},
\]
read as ``grounds''. The target statement is therefore the
well-typed atomic formula
\[
G\bigl(\operatorname{custFact}(b),\,
       \operatorname{nonjustFact}(b,c)\bigr).
\tag{G}
\]
(O$_3$) records the standing requirement of contemporary grounding
theories that grounding implies truth (or obtainment) of the
relata. (O$_1$) and (O$_2$) pin down what it is for the reified
fact-objects to obtain; without them, $\operatorname{custFact}$ and
$\operatorname{nonjustFact}$ would be ``floating'' objects that
exist independently of whether the base condition holds, which is
ontologically profligate. The proof of Proposition~\ref{prop:underdet}
does not depend on (O$_1$)---(O$_3$); they are stated here so that
the reification device is not a free-floating formalism.

\subsection{The Class of Admissible Reconstructions $\mathcal K$}
\label{sec:admissible}

For the underdetermination result to be a substantive claim about
\emph{this} encoding rather than a triviality, the class of
admissible reconstructions must be specified. Let
\[
L^+ \;=\; L_0 \,\cup\,
\{\,\operatorname{custFact},\,
  \operatorname{nonjustFact},\,
  \operatorname{Obtains},\,
  G \,\}.
\]
A structure $\mathcal M^+$ for $L^+$ is \emph{admissible} (i.e.\
$\mathcal M^+\in\mathcal K$) iff it satisfies the following
constraints:
\begin{enumerate}
\item $\mathcal M^+$ satisfies the mechanism axioms (M$_{0a}$) and
      (M$_{0b}$) on its $L_0$-reduct;
\item the typing of $\operatorname{custFact}$,
      $\operatorname{nonjustFact}$, $G$ is respected: the function
      symbols land in $\mathit{Fact}$, and $G$ is a binary relation
      on $\mathit{Fact}\times\mathit{Fact}$;
\item the structural axioms (O$_1$)--~(O$_3$) are satisfied;
\item \emph{no further bridge axiom} connects
      $\operatorname{Custom}$, $\operatorname{Just}$, or
      $\operatorname{Bel}$ to the extension of $G$ beyond (O$_1$)
      and (O$_2$). In particular, no axiom of the form
      $\varphi(b,c)\rightarrow G(\operatorname{custFact}(b),
      \operatorname{nonjustFact}(b,c))$, with $\varphi$ an $L_0$
      formula, is imposed.
\end{enumerate}
The fourth condition is the operative one. It is what makes the
underdetermination result interesting: the encoding does not
\emph{by stipulation} settle the question we are trying to put to
it. Other choices of $\mathcal K$ are possible; the result is
relative to the one stated.

\subsection{Proposition~2: Implicit Definability Failure over $\mathcal K$}
\label{sec:prop2}

\begin{lemma}[$L_0$-reduct invariance]
\label{lem:reduct-invariance}
Let $\mathcal M_1^+,\mathcal M_2^+\in\mathcal K$ with identical
$L_0$-reducts:
$\mathcal M_1^+\!\upharpoonright L_0
 =\mathcal M_2^+\!\upharpoonright L_0$.
Then, for every $L_0$-sentence $\varphi$,
$\mathcal M_1^+\models\varphi \Leftrightarrow
 \mathcal M_2^+\models\varphi$.
\end{lemma}

\begin{proof}
Induction on the complexity of $L_0$ formulas. Atomic $L_0$ formulas
have the same truth value in the two structures because their
interpretations are identical on the common $L_0$-reduct. The
Boolean cases follow immediately. The quantified cases follow because
the relevant domains and interpretations of all $L_0$ symbols are the
same in the two structures. $G$, $\operatorname{Obtains}$,
$\operatorname{custFact}$, $\operatorname{nonjustFact}$ are not in
$L_0$ and therefore cannot affect the truth of any $L_0$ sentence.
\end{proof}

\begin{proposition}[Implicit definability failure of the grounding atom]
\label{prop:underdet}
There exist admissible structures $\mathcal M_1^+,\mathcal M_2^+\in\mathcal K$
with identical $L_0$-reducts such that
\[
\mathcal M_1^+\models
G\bigl(\operatorname{custFact}(b),\operatorname{nonjustFact}(b,c)\bigr),
\qquad
\mathcal M_2^+\not\models
G\bigl(\operatorname{custFact}(b),\operatorname{nonjustFact}(b,c)\bigr).
\]
Consequently, the grounding atom (G) is \emph{not implicitly
definable} in $L_0$ over $\mathcal K$.
\end{proposition}

\begin{proof}
Take a fixed admissible $L_0$-structure $\mathcal M$ and construct
\[
\mathcal M_1^+=\langle\mathcal M,\,G_1\rangle,\qquad
\mathcal M_2^+=\langle\mathcal M,\,G_2\rangle,
\]
where $G_1$ contains
$\bigl(\operatorname{custFact}(b),\operatorname{nonjustFact}(b,c)\bigr)$
and $G_2$ does not. Choose $\mathcal M$ so that
$\operatorname{Custom}(b)$ and $\neg\operatorname{Just}(b,c)$ hold; then by
(O$_1$)--(O$_2$) both relata of the target atom obtain, and (O$_3$) holds in
$\mathcal M_1^+$ (and vacuously in $\mathcal M_2^+$). Hence both structures
are admissible, i.e.\ $\mathcal M_1^+,\mathcal M_2^+\in\mathcal K$. By
construction,
$\mathcal M_1^+\!\upharpoonright L_0
 =\mathcal M_2^+\!\upharpoonright L_0$.
By Lemma~\ref{lem:reduct-invariance}, the two structures agree on
every $L_0$-sentence. Nevertheless, they disagree on (G). Hence no
$L_0$-sentence determines whether (G) holds. This is exactly the
failure of implicit definability: the $L_0$-reduct alone does not
fix the extension of $G$ at the target pair.
\end{proof}

\begin{proposition}[Definitional irreducibility of the grounding atom]
\label{prop:definability}
Let $\mathcal M_1^+,\mathcal M_2^+$ be the pair of
Proposition~\ref{prop:underdet}. There is no $L_0$-formula
$\theta(b,c)$ with free variables of sort $B$ and $\mathit{Cont}$ such
that, for both $i=1,2$,
\[
\mathcal M_i^+\models
\forall b\,\forall c\,
\bigl[
\theta(b,c)\leftrightarrow
G(\operatorname{custFact}(b),\operatorname{nonjustFact}(b,c))
\bigr].
\]
\end{proposition}

\begin{proof}
The atom (G) is true in $\mathcal M_1^+$ and false in
$\mathcal M_2^+$, while $\theta$'s interpretation depends only on
the $L_0$-reduct, which is identical in both. Hence $\theta$ has
the same truth value at $(b,c)$ in both structures, and the
biconditional cannot hold in both simultaneously.
\end{proof}

\begin{proposition}[Stability under arbitrary $L_0$-theories]
\label{prop:stability}
Let $T$ be any set of $L_0$-sentences and let
$\mathcal M\models T$. Then the enrichment pair of
Proposition~\ref{prop:underdet} satisfies
$\mathcal M_1^+\models T$ and $\mathcal M_2^+\models T$. Consequently,
no $L_0$-theory can separate the two grounding readings while
preserving the reduct $\mathcal M$; any constraint on $G$ must be
stated in the enriched signature itself.
\end{proposition}

\begin{proof}
The two enrichments share the same $L_0$-reduct, so by
Lemma~\ref{lem:reduct-invariance} they agree on every
$L_0$-sentence; in particular, on every sentence of $T$.
\end{proof}

The three propositions jointly establish the underdetermination
claim at three different levels of strength. The mechanism that
produces them is general: it would also produce them for a relation
$G$ between arbitrary other reified objects, and the present result
is not a \emph{content-specific} limit of the encoding but a
consequence of the simple fact that $G$ is not in $L_0$. The
\emph{interest} of the result is not the mechanism but the
\emph{content} of the symbol on which it acts: $G$ is the one
relation that the contemporary grounding literature standardly
deploys to regiment the kind of explanatory claim at issue in
Hume, and the result shows that regimenting that claim requires
semantic resources beyond the $L_0$-reduct.

\subsubsection{Definability qualification: explicit, implicit, and the Beth anchor}
\label{sec:definability-qual}

The underdetermination result of
\S\ref{sec:enrichment}--\ref{sec:prop2} is best read against the
classical background of definability theory.

\begin{definition}[Explicit definability]
A relation $R$ on a structure $\mathcal M$ is \emph{explicitly
definable in $L_0$} iff there is an $L_0$-formula $\theta(\bar x)$
with $\mathcal M\models\forall\bar x\,(R(\bar x)\leftrightarrow\theta(\bar x))$.
\end{definition}

\begin{definition}[Implicit definability]
A relation $R$ is \emph{implicitly definable in $L_0$ over a class
$\mathcal K$} iff for every $\mathcal M_1,\mathcal M_2\in\mathcal K$,
$\mathcal M_1\!\upharpoonright L_0 = \mathcal M_2\!\upharpoonright L_0$
implies $\mathcal M_1\!\upharpoonright R = \mathcal M_2\!\upharpoonright R$.
\end{definition}

Proposition~\ref{prop:definability} is a failure of explicit
definability; Proposition~\ref{prop:underdet} is a failure of
implicit definability; Proposition~\ref{prop:stability} says the
implicit failure persists under any $L_0$-theory. The
\emph{Beth definability theorem} (Beth 1953) guarantees, under
suitable conditions, that implicit definability entails explicit
definability. The conditions of Beth's theorem are
\emph{not} satisfied in the present setting: the class
$\mathcal K$ is closed under $L_0$-elementary equivalence between
its members but not under the constructions Beth's proof requires
(arbitrary ultrapowers or unions of chains of $L_0$-elementary
substructures). The present argument therefore does \emph{not} rely
on Beth's theorem; it is a direct construction of a pair of
admissible structures with the required properties. The theorem is
mentioned here only to fix vocabulary and to make the
underdetermination claim comparable to the standard model-theoretic
literature.

\subsection{What Richer Languages Add---and Do Not Add}
\label{sec:richer-languages}

The negative result concerns the chosen extensional signature. Any
richer representation must therefore add semantic resources rather
than merely re-interpret the original $L_0$ vocabulary. We record
three such resources; the choice among them is the philosophical
work, and this section is deliberately neutral among them.

\subsubsection{Modal enrichment $\Box_s$}
\label{sec:modal-enrichment}

A modal enrichment introduces a unary operator $\Box_s$ together
with an accessibility relation $R$ on the (extended) domain. The
subscript $s$ is an explicit \emph{index} that fixes the intended
reading: $s$ may be (i) a source-reliability index (the impression's
provenance), (ii) an epistemic accessibility index, or (iii) a
metaphysical necessity index. The choice is interpretive; for
example, the Stoic clause ``could not arise from what is not'' is
naturally read as a \emph{source-reliability} claim and would be
regimented as $\Box_{\mathrm{src}}\,\varphi$, while a metaphysical
necessity reading of the same clause would be $\Box_{\mathrm{met}}\,\varphi$.
Conflating these readings is a common error, and we do not pre-empt
it here. Formulas such as $\Box_s\varphi$ admit a standard
translation into FOL over enriched structures (the standard
translation of modal logic into first-order logic with a binary
accessibility relation); the translation changes the signature, and
even then it captures truth in a model, not the hyperintensional
``because'' (cf.\ \S\ref{sec:enrichment}). Modal enrichment is
\emph{necessary but not sufficient} for representing the
explanatory residue.

\subsubsection{Justification-logic enrichment}
\label{sec:just-enrichment}

A justification-logic enrichment introduces explicit evidence terms
of sort $E$ and expressions of the form
\[
t:\varphi,
\]
read as ``$t$ is a justification for $\varphi$''. The construction
$t:\varphi$ is \emph{not} an $L_0$ formula: it is a separate
syntactic category, governed by a separate proof system, in which
``$t:\varphi$'' is a judgement rather than a proposition. In
particular, the term language has its own formation rules, and the
factivity schema $t:\varphi\rightarrow\varphi$ (or, in weaker
variants, $t:\varphi\rightarrow(\varphi\lor\neg\varphi)$) is an
inference rule of the metalanguage, not a $L_0$ axiom. Conflating
$t:\varphi$ with a $L_0$ predicate ``$t$ justifies $\varphi$'' is a
frequent error and obscures the way in which evidence becomes
representable. With this enrichment, the Humean situation is
represented structurally: for a causal belief $a$, the theory
contains $c:\operatorname{Custom}(a)$ (the custom term attests the
production claim) while, at the meta-level, the theory has no
theorem of the form $r:\operatorname{Just}(a,c)$ for any closed
rational-ground term $r$. The two roles that the extensional
encoding collapses---``what produced the belief'' and ``what
justifies it''---are separated at the level of the term language.
This is the minimal positive moral of our negative result: the
distinction requires a language in which \emph{evidence} is an
explicit first-class object.

\subsubsection{Grounding enrichment}
\label{sec:grounding-enrichment}

A grounding enrichment introduces a partial-grounding connective
$<$ on $\mathit{Fact}\times\mathit{Fact}$ with the standard
regimentations: irreflexivity, asymmetry and transitivity
(Fine 2012; Correia \& Schnieder 2012). Under the reification
device of \S\ref{sec:enrichment}, the Humean exclusion is written
$G(\operatorname{custFact}(b),\operatorname{nonjustFact}(b,c))$ as a
grounding claim, with $G$ playing the role of $<$ in a regimentation
that respects (O$_3$). The direction of determination---the feature
absent from model theory---is the very content of the connective.
Note that grounding is hyperintensional, so the failure of $L_0$ to
fix the extension of $G$ is not specific to FOL's extensionality;
it persists for modal extensions as well. This is why modal
enrichment, while necessary, is not sufficient.

\subsubsection{Disjunctive conclusion}
\label{sec:disjunctive}

We draw no conclusion about which of the three enrichments is
``the'' correct one. The conclusion is disjunctive: \emph{some}
hyperintensional, modal, or explicit-evidence resource is required,
and the $L_0$-reduct is not such a resource at this granularity.
The result is therefore about this encoding, at this level of
granularity, under extensional semantics. It is not a universal
claim about the expressive limits of first-order logic, nor a
denial that any richer first-order language could represent the
relevant distinction.

\subsection{Historical Anchoring: The Four-Level Humean Analysis}
\label{sec:hume-historical}

The Humean readings of \S\ref{sec:hume-readings} are not
historically neutral. To keep the formal result from silently
importing a contested interpretive thesis, we record a four-level
decomposition of the claim and tag each level with its textual
anchoring and its standing in the literature.

\begin{center}
\begin{tabular}{p{0.18\linewidth} p{0.42\linewidth} p{0.30\linewidth}}
\hline
\textbf{Level} & \textbf{Claim} & \textbf{Status} \\
\hline
Psychological-genetic &
Causal expectations are formed by custom/habit, not by demonstrative
reason. &
Directly defensible from Hume's text (T 1.3.14.20, SBN 164--65;
E 7.6, SBN 63; T 1.3.16.9, SBN 178--9). \\
Epistemic &
These beliefs are not demonstratively justified. &
Defensible, but the term ``justification'' must be unpacked
(§\ref{sec:hume-readings}). \\
Normative exclusion &
Custom-production excludes rational justification. &
Contested contemporary reconstruction; alternative readings exist
(H-II, naturalistic; Della Rocca 2010 on PSR as demotion vs\
rejection; Pruss 2006 on Hume as evaporationist). \\
Grounding &
The custom-production relation \emph{grounds} the absence of
justification. &
Contemporary regimentation in the sense of Fine (2012) and
Schnieder (2011); not directly attributable to Hume. \\
\hline
\end{tabular}
\end{center}

The formal apparatus above operates at levels 2--4 only with the
qualifications noted. We do not attribute a contemporary grounding
relation, nor a fully articulated theory of epistemic justification,
to Hume. We use grounding only to regiment one strong contemporary
reconstruction of the explanatory force that an interpreter may
take Hume's appeal to custom to possess; alternative reconstructions
remain live, and the formal apparatus is invariant under the choice
between H-I and H-II modulo the annotations recorded in
\S\ref{sec:prop1}--\S\ref{sec:prop2}.

\subsection{The Stoic Modal Clause: A Note}
\label{sec:stoic-modal-note}

The ancient definition of the kataleptic impression (DL 7.46
[Dorandi 2013]; M 7.248, 7.402 [Bury 1935]) contains a literal modal
clause: the impression is ``of such a kind as could not arise from
what is not''. The minimal dependency skeleton (S) of
\S\ref{sec:stoic-skeleton} \emph{does not} contain this clause.
Encoding it requires either (i) a modal operator $\Box_s$ with an
appropriately indexed reading (§\ref{sec:modal-enrichment}), or
(ii) a separate monadic predicate $\operatorname{Ver}(i)$ on
impressions together with an explicit axiom
$\operatorname{Kat}(i)\rightarrow\operatorname{Ver}(i)$. The choice
between (i) and (ii) is interpretive. We do not pre-empt it: the
present formalization is at the level of the minimal dependency
skeleton, and the modal clause is identified as a separate residue
to be added by whichever enrichment the reader finds defensible.

The separation of (S) from the modal clause is methodologically
important. It is the formal counterpart of the historical point
that katalepsis and \emph{epist\=em\=e} are
\emph{not} identical in Stoic epistemology: katalepsis is the
cognition of a kataleptic impression, while \emph{epist\=em\=e}
is a stable, secure cognitive state distinct from mere cognitive
uptake. The minimal dependency skeleton leaves room for this
distinction; encoding the Stoic position more fully would require
additional predicates and a separate argument that we do not
undertake here.

\subsection{Verification Status and Reproducibility}
\label{sec:verification}

The formal core has been checked at two levels.

\emph{First}, the finite Boolean fragment underlying
Proposition~\ref{prop:strength} and
Proposition~\ref{prop:bridge-collapse} was exhaustively enumerated
for the single-$B$, single-$\mathit{Cont}$ case. Sixteen interpretations were
generated. The consequence $T_2\land M_0\models T_1$ was confirmed
(zero falsifying interpretations); the stated countermodel to
$T_1\land M_0\models T_2$ was verified; the bridge collapse
$T_1\land B_0\models T_2$ was verified (zero falsifying
interpretations in the Boolean fragment); and the characterization
$(\star)$ of the separation between $T_1$ and $T_2$ over $M_0$ was
confirmed. The full table of 16 interpretations is given in
\S\ref{sec:verification-table} below.

\emph{Second}, the model pair required for
Proposition~\ref{prop:underdet} was explicitly constructed with
identical $L_0$ interpretations and distinct $G$ extensions. This
machine check does not replace the general semantic proof: the
latter is supplied by Lemma~\ref{lem:reduct-invariance}. The
computational check is therefore a reproducibility and
error-detection layer, while the lemma-and-proposition argument
remains the mathematical proof presented in this section.

The associated reproducibility materials implement exactly this
finite check and the two-model construction. They should be treated
as verification artefacts, not as substitutes for the formal
argument in the main text. Both the Python script
\texttt{core\_formal\_model\_check.py} and the documented
run-output are shipped with the paper; the script's results are
stable across the Python versions tested (3.10, 3.11, 3.12) and
are independent of any third-party library.

\subsubsection{Table A: Sixteen interpretations of the Boolean core}
\label{sec:verification-table}

The Boolean fragment consists of the four predicates
$\operatorname{Causal}$, $\operatorname{Custom}$, $\operatorname{Bel}$,
$\operatorname{Just}$ evaluated on a single $(b,c)$ pair. There are
$2^4=16$ interpretations. The columns are: truth values of the
four predicates; truth of $M_0$ (i.e.\ M$_{0a}\land$M$_{0b}$); truth
of $T_1$; truth of $T_2$. ``\checkmark'' denotes true; ``$\times$''
denotes false.

\begin{center}
\small
\begin{tabular}{cccc|ccc|c}
\hline
$\operatorname{Causal}$ & $\operatorname{Custom}$ &
$\operatorname{Bel}$ & $\operatorname{Just}$ & $M_0$ & $T_1$ & $T_2$ & note \\
\hline
0 & 0 & 0 & 0 & \checkmark & \checkmark & \checkmark & \\
0 & 0 & 0 & 1 & \checkmark & \checkmark & \checkmark & \\
0 & 0 & 1 & 0 & \checkmark & \checkmark & \checkmark & \\
0 & 0 & 1 & 1 & \checkmark & \checkmark & \checkmark & \\
0 & 1 & 0 & 0 & \checkmark & \checkmark & \checkmark & \\
0 & 1 & 0 & 1 & \checkmark & \checkmark & \checkmark & \\
0 & 1 & 1 & 0 & \checkmark & \checkmark & \checkmark & \\
0 & 1 & 1 & 1 & \checkmark & \checkmark & $\times$ & \textbf{unique countermodel} \\
1 & 0 & 0 & 0 & $\times$ & \checkmark & \checkmark & M$_0$a \& M$_0$b violation \\
1 & 1 & 0 & 0 & $\times$ & \checkmark & \checkmark & M$_0$b violation \\
1 & 0 & 1 & 0 & $\times$ & \checkmark & \checkmark & M$_0$a violation \\
1 & 0 & 0 & 1 & $\times$ & $\times$ & \checkmark & \\
1 & 0 & 1 & 1 & $\times$ & $\times$ & \checkmark & \\
1 & 1 & 0 & 1 & $\times$ & $\times$ & \checkmark & \\
1 & 1 & 1 & 0 & \checkmark & \checkmark & \checkmark & \\
1 & 1 & 1 & 1 & \checkmark & $\times$ & $\times$ & \\
\hline
\end{tabular}
\end{center}

Two facts are worth highlighting.
\begin{enumerate}
\item The countermodel to $T_1\land M_0\models T_2$ is
\emph{unique}: of the 16 interpretations, exactly one
($(0,1,1,1)$) falsifies the entailment. The finite Boolean
fragment is therefore as informative as the general
Lemma~\ref{lem:reduct-invariance} for the single-pair case.
\item The bridge collapse of
Proposition~\ref{prop:bridge-collapse} is sharp: adding $B_0$
eliminates \emph{all} falsifying interpretations of
$T_1\land B_0\land M_0\models T_2$ (verified by adding the
$B_0$ column to the same 16-row table).
\end{enumerate}
The computational evidence is therefore a finite-validation
artefact, not a substitute for the general proof, and is consistent
with all three propositions of \S\ref{sec:prop2} as well as with the
characterization $(\star)$ of
\S\ref{sec:bridge-collapse}.

\subsection{Scope of the Formal Result}
\label{sec:formal-scope}

The formal result is relative to the language, ontology, and
interpretive choices specified above. It does not establish that:
\begin{enumerate}
\item no richer first-order language could represent the relevant
      distinction;
\item modal logic is philosophically superior to grounding or
      justification logic;
\item Stoic katalepsis is identical to any contemporary theory of
      reliabilism;
\item Hume denies every form of epistemic justification, or that
      the ``strong exclusion reconstruction'' (H-I) is the uniquely
      correct reading of Hume's appeal to custom;
\item the formal reconstruction uniquely determines the historical
      interpretation of the Stoic or Humean texts.
\end{enumerate}
The claim is narrower: an extensional reconstruction at the
specified level of granularity, on the specified class of
admissible enrichments $\mathcal K$, can preserve material content
while leaving the explanatory and modal relations semantically
underdetermined.

\begin{quote}
This article does not show that first-order logic cannot represent
grounding, modality, or epistemic normativity. It shows that, for
the explicitly specified extensional signature $L_0$ and class of
admissible reconstructions $\mathcal K$, the intended explanatory
relation between custom-production and the absence of rational
warrant is not fixed by the $L_0$-reduct; representing it requires
additional structure whose philosophical interpretation remains
independently contestable.
\end{quote}
